\chapter{Examples of oriented images and weak derivatives}\label{sec:examples-current-images}
The signed multiplicity formula in \cite{Sim18}, Chapter 6, pp.~174--175, distinguishes an oriented current from the measure of its parametrized image. These examples use the area calculations in Section \ref{sec:examples-area-critical} and the local Lipschitz pushforward theorem.

\begin{example}[Folding a rectangle cancels its oriented image]\label{ex:current-folded-rectangle}
Let $R=(-1,1)\times(0,1)$ carry the standard orientation and put $T=\llbracket R\rrbracket$ in $\mathbb R^2$. For $f(x,y)=(|x|,y)$,
\[
 f_\#T=0,\qquad f_\#(\partial T)=0,
 \qquad f_\#(\|T\|)=2\mathcal L^2\llcorner((0,1)^2).
\]
Here the last pushforward is the ordinary pushforward of the area measure; for this map its area Jacobian is one almost everywhere, so it is also the unsigned area image. The current and its measure have different image operations.
\end{example}
\begin{proof}
Away from $x=0$, the derivative is $\operatorname{diag}(\operatorname{sign}x,1)$, with signed determinant $+1$ on the right half and $-1$ on the left. Every point in $(0,1)^2$ has two preimages with these opposite signs and multiplicity one. The critical/nondifferentiability line has zero area and contributes nothing to the tangential integral. The signed multiplicity theorem therefore gives zero, while substitution on the two halves gives measure density two. The rectangle current is integral and compactly supported, and $f$ is Lipschitz, so Theorem \ref{thm:lipschitz-current-pushforward} gives $f_\#\partial T=\partial(f_\#T)=0$. The vertical left and right boundary edges cancel after folding; the horizontal edges are retraced in opposite directions. No nonzero boundary current survives at the fold line.
\end{proof}

\begin{example}[A power map preserves orientation and multiplies disk multiplicity]\label{ex:current-complex-power-disk}
Let $q\geq1$ be an integer, $R>0$, and let $T=\llbracket D_R\rrbracket$ be the standard oriented complex disk as a two-current in $\mathbb R^2$. For $f(z)=z^q$,
\[
 f_\#T=q\llbracket D_{R^q}\rrbracket,
 \qquad f_\#(\partial T)=q\partial\llbracket D_{R^q}\rrbracket,
 \qquad \mathbf M(f_\#T)=\pi qR^{2q}.
\]
The map is locally Lipschitz and proper on the compact support of $T$; it need not be globally Lipschitz when $q>1$.
\end{example}
\begin{proof}
Example \ref{ex:area-complex-power-map} computes the Jacobian and the $q$ distinct roots of each nonzero image point. Multiplication by the nonzero complex derivative preserves real orientation, since its real determinant is its squared modulus. Thus all $q$ sheets have positive sign, and Theorem \ref{thm:integer-rectifiable-geometric-image}, applied after a Lipschitz extension near the compact disk, gives image multiplicity $q$. The origin is an exceptional null image point; the branch point contributes no separate mass. The disk has mass $\pi R^{2q}$ after changing radius, giving the displayed mass. Boundary commutation in Theorem \ref{thm:lipschitz-current-pushforward} proves the circle-current identity with its induced counterclockwise orientation. This separates the geometric image disk, whose set has multiplicity one, from its pushed current.
\end{proof}

\begin{example}[Uniformly small oscillations retain nonconvergent derivatives]\label{ex:weak-minors-oscillating-maps}
Fix $0<a,b<1$ and, on $V=(0,2\pi)^2$, set
\[
 u_j(x,y)=\left(x+\frac aj\sin(jx),\ y+\frac bj\sin(jy)\right),
 \qquad j=1,2,\ldots.
\]
Then $u_j$ converge uniformly to the identity and have a common Lipschitz bound. Their Jacobians are
\[
 \det Du_j=(1+a\cos(jx))(1+b\cos(jy)).
\]
They converge weakly against $L^1(V)$ tests to one, although
\[
 \int_V|\partial_xu_j^1-1|\,dx\,dy=8\pi a
\]
for every $j$. Thus the pullback convergence argument cannot be replaced by strong convergence of derivatives. Extending the displayed maps to the whole plane, they also satisfy $(u_j)_\#\llbracket V\rrbracket=\llbracket V\rrbracket$ for every $j$.
\end{example}
\begin{proof}
The map error is at most $\sqrt{a^2+b^2}/j$. Its derivative is diagonal with entries in $[1-a,1+a]$ and $[1-b,1+b]$, giving a common operator bound $1+\max\{a,b\}$ and the Jacobian formula. Proposition \ref{prop:uniform-lipschitz-weak-minors} gives convergence against compactly supported $L^1$ tests. The Jacobians are uniformly bounded, so truncating an arbitrary $L^1(V)$ test away from the boundary extends the conclusion to all such tests. Finally, $\partial_xu_j^1-1=a\cos(jx)$ and $\int_0^{2\pi}|\cos(jx)|\,dx=4$ for integer $j$. Multiplication by the length $2\pi$ of the other interval gives $8\pi a$. The nonvanishing norm disproves strong $L^1$ convergence; the example does not assert pointwise divergence at every coordinate. Each coordinate map is strictly increasing because its derivative is positive, and it fixes the endpoints zero and $2\pi$. Hence $u_j$ maps $V$ diffeomorphically onto itself with positive orientation. Theorem \ref{thm:integer-rectifiable-geometric-image} then gives the exact current identity, since each image point has one positively oriented preimage.
\end{proof}

\begin{example}[An injective map distinguishes pushed length from pushed current mass]\label{ex:current-fat-cantor-measure-distinction}
Let $f:[0,1]\to[0,1/2]$ be the strictly increasing Lipschitz map of Example \ref{ex:area-fat-cantor-critical-map}, extended constantly to the whole line, and let $T=\llbracket[0,1]\rrbracket$ with its positive orientation. Then
\[
 f_\#T=\llbracket[0,1/2]\rrbracket,
 \qquad \|f_\#T\|=\mathcal L^1\llcorner[0,1/2],
 \qquad f_\#\|T\|=\mathcal L^1\llcorner[0,1/2]+\nu,
\]
where $\nu$ is the nonzero singular continuous measure of mass $1/2$ from that example. There is no cancellation: the difference comes from the area Jacobian in the current formula. Also $f_\#\partial T=\delta_{1/2}-\delta_0$.
\end{example}
\begin{proof}
For a smooth compactly supported test one-form $\phi(y)\,dy$, the geometric pushforward is
\[
 (f_\#T)(\phi(y)\,dy)=\int_0^1\phi(f(x))f'(x)\,dx
 =\int_0^{1/2}\phi(y)\,dy.
\]
The last equality follows from the injective weighted area formula; $f'\geq0$ almost everywhere. This gives the current and its mass measure. By contrast the ordinary pushforward of length tests against $\phi(f(x))$ without the derivative, and Example \ref{ex:area-fat-cantor-critical-map} computes its singular summand $\nu$. On the positive-measure set $E$, the derivative is zero almost everywhere, so that part contributes zero to the current image and mass $1/2$ to the ordinary pushed length. Positive orientation at the regular points rules out sign cancellation. Finally $\partial T=\delta_1-\delta_0$ and $f(1)=1/2$, $f(0)=0$ give the boundary formula, consistently with boundary commutation. The unsigned area image uses the same Jacobian as the current and therefore has mass $1/2$, while the ordinary pushed length retains total mass one.
\end{proof}

\begin{example}[Mass convergence can have an infinite-length union of carriers]\label{ex:mass-convergence-infinite-carrier-union}
In $\mathbb R^2$, let $I$ be the rightward current of $[0,1]\times\{0\}$ and let $S_j$ be the rightward current of $[0,1/j]\times\{1/j\}$, for integers $j\geq1$. Put $T_j=I+S_j$. These integral currents have a common compact support container and
\[
 \mathbf M(T_j-I)=\frac1j\longrightarrow0,
 \qquad \mathbf M(\partial T_j-\partial I)=2.
\]
Their carrier union has infinite length in the compact square $[0,1]^2$. The boundaries converge weakly to $\partial I$ but are not Cauchy in mass.
\end{example}
\begin{proof}
The segments $S_j$ have length $1/j$ and lie at distinct positive heights, disjoint from $I$. Thus the first mass equality is exact. Countable additivity gives length $1+\sum_{j\geq1}1/j=\infty$ for their union. For every $r>0$, all segments with $j>\sqrt2/r$ lie in $B_r(0)$ and their lengths still have a divergent sum. Thus the union is locally infinite near the origin, as permitted in Lemma \ref{lem:common-countable-current-carrier}.

The extra boundary is
\[
 \partial S_j=\delta_{(1/j,1/j)}-\delta_{(0,1/j)}.
\]
Its two points are distinct, so its mass is two. Both points tend to the origin; testing against a smooth compactly supported scalar function shows that their signed difference tends to zero. For $j\ne k$, the four endpoints of $S_j$ and $S_k$ are distinct, so $\mathbf M(\partial T_j-\partial T_k)=4$. Hence the boundaries are not mass-Cauchy, despite having the uniform bound $\mathbf M(\partial T_j)=4$. This explains why Corollary \ref{cor:integral-current-local-normal-completeness} assumes boundary mass convergence separately. The weak compactness theorem is not contradicted: here the weak boundary limit exists and the stronger mass convergence fails.
\end{proof}
